\section{Spatiotemporal predictions from the model}\label{sec:spatiotemporal_predictions}

Throughout this chapter, we present a novel physics-inspired analysis of a standard gaze-tracking experiment during visual search which yields new results and predictions of gaze dynamics. In particular, the approximate timescale of saccades has been known since the beginning of the field of eye movement research, whereas the significance of the rate-corrected local rhythmicity is a very recent result only discovered in the year of this present work. OMT as a phenomenon has been found in the literature \cite{graham_ocular_2023}, however no consensus exists as to it's importance in the functioning of gaze, and it has never been linked to a computational role of gamma-scale rhythms. Lastly, the mathematical analysis of search superdiffusion of gaze has only been recently undertaken in the literature, and without the well-motivated link to a model of gaze we present here. However, another benefit of comparing the generative unifying model presented in \cref{chap:a_unifying_model} to the data is that further dynamical predictions outside of the present literature can be made. Thus, we present experimental paradigms and features of this model which are being developed to better understand visual searching.
\subsection{Saliency-frequency coupling in contrast-dependant search tasks}
OMT frequency is here hypothesised to be the result of a harmonic oscillations in the control of the path of the eye during fixation tasks. An experiment may be designed in which targets with various degrees of contrast are presented to be searched. High contrast, localised points, such as tight, high spatial frequency Gabor patches, would produce localised saliency wells with low $\sigma_{\text{fix}}$. Whereas, low contrast, large low spatial frequency patches would produce saliency wells with a large width $\sigma_{\text{fix}}$. High contrast wells would correspond to shorter theta timescales, quadratic in the effective saliency well size as in \cref{eq:theta_rhythm}, and a microtremor rhythm which is approximately inverse to the effective saliency well size as in \cref{eq:resonant_frequency}. The effect of visual contrast has been found to be significant in gamma frequency in the visual system, and this quantitative relationship provides an experimental prediction of the relationship between three data points: $t_{\theta}, f_{\text{peak}},$ and $k_{\text{Gabor}}$. A simple regression analysis can be achieved here to find correlations between spatial frequency $k_{\text{Gabor}}$ and the visual theta-rhythm $t_{\theta}$ and the $f_{\text{peak}}$. It can then also be tested if these effects are consistent with a potential well of variable size, such that $t_{\theta} \propto \sigma_{\text{fix}}^{-2}$ and $f \propto \sigma_{\text{fix}}^{-1}$
\subsection{Smooth pursuit predictions and parametric-fitting techniques}
Another experimental paradigm which provides insight into the behaviour of the control of visual trajectories. In particular, the exact same model can be examined in the behavioural paradigm of smooth tracking of a moving object. In particular, a typical experiment would be the tracking of a smoothly moving pendulum. It is realistically found that a typical eye movement trace contains discontinuous micro-saccades during smooth pursuit, and discontinuous catch-up saccades. In particular, we expect from the model that the instability will be most present in the periods of stability when kernel deposition is strongest. Further experiments could also find the behaviour of dynamics in phase space $\alpha/\beta/\gamma/t_{\theta}$ for different neurological conditions. For instance, it is well-known that schizophrenia effects the 